
% ===========================================================================
% 9. MODELE DE DONNEES
% ===========================================================================
\section{Mod\`ele de donn\'ees et stockage}

\subsection{Indices Elasticsearch}

\begin{center}
\small
\begin{longtable}{@{}p{0.34\linewidth} p{0.18\linewidth} p{0.42\linewidth}@{}}
\toprule
\textbf{Mod\`ele d'indice} & \textbf{Source} & \textbf{Contenu} \\
\midrule
\endhead
\bottomrule
\endlastfoot
\texttt{vulnerabilities-AAAA.MM.JJ} & Logstash & R\'esultats de scan avec scores de risque \\
\texttt{filebeat-8.11.0-AAAA.MM.JJ} & Filebeat & Journaux Docker, syst\`eme, ICMP \\
\texttt{packetbeat-8.11.0-AAAA.MM.JJ} & Packetbeat & M\'etadonn\'ees de trafic r\'eseau \\
\texttt{auditbeat-8.11.0-AAAA.MM.JJ} & Auditbeat & \'Ev\'enements processus/fichiers/utilisateurs \\
\texttt{metricbeat-8.11.0-AAAA.MM.JJ} & Metricbeat & M\'etriques syst\`eme, Docker, services \\
\texttt{heartbeat-8.11.0-AAAA.MM.JJ} & Heartbeat & V\'erifications de disponibilit\'e \\
\texttt{glpi-AAAA.MM.JJ} & Logstash & \'Etat des tickets GLPI synchronis\'e \\
\bottomrule
\end{longtable}
\end{center}

\subsection{Vues de donn\'ees Kibana}

\begin{center}
\small
\begin{tabularx}{\linewidth}{@{}l X l@{}}
\toprule
\textbf{ID} & \textbf{Mod\`ele} & \textbf{Nom} \\
\midrule
\texttt{5ec5e7c3-...} & \texttt{filebeat-*} & VOC Filebeat Logs \\
\texttt{8cd32128-...} & \texttt{packetbeat-*} & VOC Network Traffic \\
\texttt{943fb14e-...} & \texttt{auditbeat-*} & VOC Audit Trail \\
\texttt{metricbeat-voc} & \texttt{metricbeat-*} & VOC Metrics \\
\texttt{heartbeat-voc} & \texttt{heartbeat-*} & VOC Uptime \\
\texttt{997fc760-...} & \texttt{vulnerabilities-*} & VOC Vulnerabilities \\
-- & \texttt{glpi-*} & GLPI Tickets \\
\bottomrule
\end{tabularx}
\end{center}

\subsection{Exemple de document de vuln\'erabilit\'e}

Un document index\'e via Logstash contient les champs suivants :

\begin{lstlisting}[language=,caption={Document de vuln\'erabilit\'e index\'e (abr\'eg\'e)}]
{
  "@timestamp": "2026-08-15T04:57:00Z",
  "host_ip": "192.168.1.1",
  "os": "Netgear WGR614v7",
  "cve": "CVE-2017-9765",
  "cvss": 9.0,
  "risk_score": 10.0,
  "severity": "Critical",
  "port": "7547/tcp",
  "source": "voc-nmap",
  "status": "active",
  "epss_score": 0.80,
  "in_kev": true,
  "exploit_available": true,
  "vulnerability_type": "Security Vulnerability",
  "cwe_ids": ["CWE-78"],
  "attack_tactics": ["Initial Access"],
  "attack_techniques": ["T1190"],
  "checklist": [ { "id": "VAL-01", "category": "Validation", "done": false } ],
  "technical_description": "## Vulnerability Overview ..."
}
\end{lstlisting}

Le champ \texttt{os} est d\'esormais pr\'esent gr\^ace \`a la d\'etection du
syst\`eme d'exploitation, et le champ \texttt{risk\_score} alimente les
tableaux de bord et les r\`egles de tickets.

\clearpage

% ===========================================================================
% 10. SECURITE
% ===========================================================================
\section{S\'ecurit\'e de la plateforme}

\subsection{Chiffrement de Kibana}

Kibana n\'ecessite des cl\'es de chiffrement pour prot\'eger les objets
sauv\'egard\'es, les sessions et les rapports. Trois cl\'es ont \'et\'e
inject\'ees via les variables d'environnement suivantes :

\begin{center}
\small
\begin{tabularx}{\linewidth}{@{}l X@{}}
\toprule
\textbf{Variable} & \textbf{R\^ole} \\
\midrule
\texttt{XPACK\_SECURITY\_ENCRYPTIONKEY} & Chiffrement des sessions et donn\'ees s\'ecuris\'ees \\
\texttt{XPACK\_ENCRYPTEDSAVEDOBJECTS\_ENCRYPTIONKEY} & Chiffrement des objets sauv\'egard\'es \\
\texttt{XPACK\_REPORTING\_ENCRYPTIONKEY} & Chiffrement des rapports \\
\bottomrule
\end{tabularx}
\end{center}

\subsection{Authentification}

Chaque service dispose de m\'ecanismes d'authentification : utilisateur
\texttt{elastic} pour Elasticsearch/Kibana, jetons d'application et
d'utilisateur pour l'API GLPI, et cl\'e API pour MISP. Les secrets sont
stock\'es dans le fichier \texttt{.env} et inject\'es dans les conteneurs.

\subsection{Isolation des conteneurs}

La plateforme utilise des r\'eseaux et volumes Docker d\'edi\'es. Le worker
de scan ex\'ecute nmap en tant que \texttt{root} avec la capacit\'e
\texttt{NET\_RAW} pour permettre le scan SYN (\texttt{-sS}) et la d\'etection
OS (\texttt{-O}) ; cette ex\'ecution privil\'egi\'ee est limit\'ee au seul
conteneur de scan. Les limites m\'emoire sont d\'efinies par service afin
d'\'eviter qu'un conteneur n'\'epuise les ressources de l'h\^ote.

\subsection{Supervision des agents}

Fleet centralise l'enr\^olement des agents et la distribution des politiques.
L'\'etat de chaque agent (\texttt{online}, \texttt{offline}, \texttt{degraded})
est visible depuis Kibana, ce qui permet de d\'etecter rapidement les agents
qui ne transmettent plus de donn\'ees.

\subsection{Transport des donn\'ees}

En environnement de laboratoire, les certificats autog\'en\'er\'es sont
utilis\'es. Les clients (MISP, GLPI, Elasticsearch) acceptent une
v\'erification SSL d\'esactivable via des variables d'environnement
(\texttt{MISP\_VERIFY\_SSL}, \texttt{GLPI\_VERIFY\_SSL}) afin de permettre un
durcissement lors d'un passage en production.

\clearpage

% ===========================================================================
% 11. ETAT OPERATIONNEL
% ===========================================================================
\section{\'Etat op\'erationnel du d\'eploiement}

\subsection{Sant\'e de la plateforme}

La plateforme est enti\`erement fonctionnelle : les 20 services Docker
(18 applicatifs + serveur Fleet + t\^ache d'initialisation ponctuelle) sont
d\'eploy\'es et, pour la quasi-totalit\'e, \emph{healthy}.

\begin{center}
\small
\begin{tabularx}{\linewidth}{@{}l X@{}}
\toprule
\textbf{Indicateur} & \textbf{Valeur} \\
\midrule
Conteneurs Docker & 20 d\'eploy\'es (18 + Fleet Server + initialisation) \\
Cluster Elasticsearch & 1 n\oe ud, statut \textcolor{vocgreen}{\textbf{green}} \\
Shards & 129 primaires actifs, 0 non-assign\'es \\
Documents de vuln\'erabilit\'es (total) & 5{,}041 \\
Documents vuln\'erabilit\'es du jour (08.15) & 2{,}992 \\
H\^otes analys\'es & 4 (dont routeur \texttt{192.168.1.1}) \\
H\^otes actifs d\'ecouverts (LAN) & 8 (routeurs, WAP, postes, serveurs) \\
Tickets GLPI & 419 (statut = nouveau, avec checklists) \\
\'Ev\'enements MISP & 1{,}300+ \\
\bottomrule
\end{tabularx}
\end{center}

\subsection{Preuves de l'extension du scan}

Le passage au scan complet a \'et\'e v\'erifi\'e en production :

\begin{itemize}[leftmargin=1.6em]
  \item le scan d\'emarre \textbf{au d\'emarrage du conteneur} puis toutes
        les 6\,h (ex\'ecution v\'erifi\'ee) ;
  \item \textbf{d\'etection du syst\`eme d'exploitation} op\'erationnelle :
        ex.\ \texttt{Linux 5.0 - 5.14} pour \texttt{192.168.1.18}, et lors de
        la d\'ecouverte : routeur Netgear WGR614, point d'acc\`es TP-LINK
        TD-W8963N, poste Windows 11 et smartphones Android ;
  \item le port \texttt{7547/tcp} (TR-069, hors de l'ancienne liste de 41
        ports) a \'et\'e d\'ecouvert, preuve que l'\textbf{int\'egralit\'e des
        ports} est maintenant scann\'ee ; il a d\'eclench\'e les tickets GLPI
        \texttt{\#418} et \texttt{\#419} (critiques, h\^ote
        \texttt{192.168.1.1}) ;
  \item l'analyse diff\'erentielle du dernier cycle : \textbf{2 nouvelles} /
        307 toujours pr\'esentes / 0 r\'esolue ;
  \item deux vuln\'erabilit\'es \`a risque 10.0 ont \'et\'e index\'ees et un
        \'ev\'enement MISP a \'et\'e cr\'e\'e ;
  \item le champ \texttt{os} est bien pr\'esent dans les documents
        Elasticsearch.
\end{itemize}

\subsection{R\'epartition des donn\'ees (indice du 08.15)}

\begin{center}
\small
\begin{tabularx}{\linewidth}{@{}l r X@{}}
\toprule
\textbf{H\^ote} & \textbf{Documents} & \textbf{Observations} \\
\midrule
\texttt{192.168.1.18} & 1{,}778 & serveur LAN (port 3306/MySQL majoritaire) \\
\texttt{192.168.184.135} & 1{,}119 & h\^ote plateforme (OpenSSH, services Docker) \\
\texttt{192.168.184.136} & 93 & second h\^ote du r\'eseau SPA \\
\texttt{192.168.1.1} & 2 & routeur d\'ecouvert via le scan complet \\
\bottomrule
\end{tabularx}
\end{center}

\subsection{Ports les plus fr\'equents}

\begin{center}
\small
\begin{tabularx}{\linewidth}{@{}l r X@{}}
\toprule
\textbf{Port} & \textbf{Documents} & \textbf{Observation} \\
\midrule
\texttt{3306/tcp} & 2{,}400 & MySQL/MariaDB (h\^ote \texttt{.18}) \\
\texttt{22/tcp} & 323 & SSH (serveurs) \\
\texttt{8080/tcp} & 252 & HTTP alternatif \\
\texttt{9200/tcp} & 15 & Elasticsearch expos\'e \\
\texttt{7547/tcp} & 2 & TR-069 (routeur) -- nouvellement d\'ecouvert \\
\bottomrule
\end{tabularx}
\end{center}

\subsection{Agents Fleet}

\begin{center}
\small
\begin{tabularx}{\linewidth}{@{}l l l X@{}}
\toprule
\textbf{Agent} & \textbf{H\^ote} & \textbf{Politique} & \textbf{\'Etat} \\
\midrule
\texttt{2fbcef12-...} & \texttt{voc-fleet-server} & fleet-server-policy-fresh-001 & \textcolor{vocgreen}{online} \\
\texttt{607d470b-...} & conteneur \'elastique & linux\_hosts & \textcolor{vocgreen}{online} \\
\texttt{f7f2c317-...} & \texttt{tbini} (VM \texttt{192.168.184.138}) & linux\_hosts & \textcolor{vocorange}{degraded} \\
\bottomrule
\end{tabularx}
\end{center}

\clearpage
